\documentclass[11pt]{amsart}
\usepackage{amsmath,amssymb,amsthm}
\usepackage[margin=1.15in]{geometry}
\usepackage[hidelinks]{hyperref}
\newtheorem{theorem}{Theorem}[section]
\newtheorem{corollary}[theorem]{Corollary}
\newtheorem{lemma}[theorem]{Lemma}
\theoremstyle{remark}
\newtheorem{remark}[theorem]{Remark}
\newcommand{\doi}[1]{\href{https://doi.org/#1}{doi:\nolinkurl{#1}}}
\newcommand{\arxiv}[1]{\href{https://arxiv.org/abs/#1}{arXiv:\nolinkurl{#1}}}
\newcommand{\eps}{\varepsilon}
\newcommand{\E}{\mathbb{E}}
\newcommand{\thit}{t_{\mathrm{hit}}}
\newcommand{\rhit}{\rho_{\mathrm{hit}}}
\setlength{\emergencystretch}{2.5em}

\begin{document}

\title[Hitting-time robustness of the hypercube]{Hitting-time robustness of the hypercube}
\author{Yitzchak Shmalo}
\address{Einstein Institute of Mathematics, The Hebrew University of Jerusalem, Jerusalem, Israel}
\email{yspennstate@gmail.com}
\subjclass[2020]{60J10, 05C81, 60G42}
\keywords{random walk, hypercube, hitting time, adversarial perturbation, Markov chain robustness}

\begin{abstract}
Zuckerman defined the hitting-time robustness of a random walk as the largest $\eps$ for which an adversary who may
change every transition probability by a factor in $[1-\eps,1+\eps]$, depending on the whole history, cannot raise the
maximum expected hitting time by more than a constant factor. He asked for its value on hypercubes. For the simple
random walk on $\{0,1\}^d$ we show that the strongest such adversary with $\eps=c/d$ raises the maximum hitting time by
a factor tending to $2^c$. Hence the largest $\eps$ for which the factor stays at most $C$ is $(\log_2 C+o(1))/d$. For
the lazy walk the factor tends to $e^c$ and that largest $\eps$ is $(\ln C+o(1))/d$. Both are of order one over the
diameter, which confirms on the hypercube Zuckerman's conjecture that the robustness of every graph is at least of that
order.
\end{abstract}

\maketitle

\section{Introduction}\label{sec:intro}

Let $G$ be a connected graph and $P$ the transition matrix of the simple random walk on $G$. Following
Zuckerman~\cite{Zuc26}, a \emph{standard $\eps$-adversary} ($0<\eps<1$) moves the walk from its current vertex $v$ to a
neighbour $w$ with a probability of its choosing in $[(1-\eps)P(v,w),(1+\eps)P(v,w)]$, the probabilities from $v$
summing to one; the choice may depend on the whole history. Write $H_v$ for the hitting time of $v$, $h^+(u,v)$ for
the supremum of $\E_u H_v$ over all standard $\eps$-adversaries, $\thit^+(\eps)=\max_{u,v}h^+(u,v)$, and
$\thit=\max_{u,v}\E_uH_v$ for the unperturbed walk. The \emph{hitting-time robustness} is
\[
\rhit(C)=\sup\{\eps:\ \thit^+(\eps)\le C\,\thit\}.
\]
When $C$ is omitted, $C$ is a sufficiently large constant, as in~\cite[Definition~1.9]{Zuc26}.
Zuckerman proved $\rhit=\Omega(1/t_{\mathrm{mix}})$ for every graph~\cite[Theorem~1.10]{Zuc26}, where
$t_{\mathrm{mix}}$ is the mixing time (of the lazy walk when $G$ is bipartite~\cite[Remark~2.9]{Zuc26}), and
$\rhit=\Theta(1/\log n)$ for constant-degree expanders on $n$ vertices~\cite[Theorem~1.12]{Zuc26}. He conjectured $\rhit=\Omega(1/\Delta)$ with $\Delta$ the diameter, proved it for trees~\cite[Theorem~1.11]{Zuc26}, and listed the hypercube
among the natural graphs whose robustness is to be determined~\cite[Section~1.10]{Zuc26}.

On the hypercube $Q_d=\{0,1\}^d$, where $n=2^d$ and $\Delta=d$, these results leave a logarithmic gap. The lazy walk
on $Q_d$ mixes in time $\Theta(d\log d)$~\cite[Section~18.2]{LPW17}, so~\cite[Theorem~1.10]{Zuc26} gives
$\rhit=\Omega(1/(d\log d))$. In the other direction, the adversary constructed in the proof of~\cite[Theorem~4.5]{Zuc26},
which we call the tilted adversary, moves each step away from the target. On $Q_d$ its transition probabilities lie in
$[(1-\eps)/d,\,1/((1-\eps)d)]$, so it is a standard $\eps/(1-\eps)$-adversary, and it multiplies $\thit$ by about
$e^{\eps d/2}$. This gives $\rhit=O(1/d)$. The proof of~\cite[Theorem~1.10]{Zuc26} compares the perturbed walk with the
original over a mixing period, which on $Q_d$ has length of order $d\log d$, and the logarithm is lost there. On $Q_d$
the supremum over adversaries can be computed instead. Call \emph{extremal} the adversary that at every vertex sends the
least mass toward the target that the bounds allow. A supermartingale comparison (Lemma~\ref{lem:reduction}) shows that
the extremal adversary attains the supremum. It depends only on the distance to the target, and under it the distance is
a birth--death chain on $\{0,\dots,d\}$ with stationary weights $w_0=1,w_1,\dots,w_d$. For $\eps=c/d$, uniformly in $j$,
\begin{gather*}
\thit^+(\eps)=W\,(1+o(1)),\qquad W=\sum_{j=0}^dw_j,\\
\ln\frac{w_j}{\binom dj}=c\,F\Bigl(\frac jd\Bigr)+O\Bigl(\frac1d\Bigr),\qquad
F(x)=\begin{cases}-\ln(1-x),&0\le x\le\frac12,\\ \ln(4x),&\frac12\le x\le1\end{cases}
\end{gather*}
(Lemmas~\ref{lem:reduction}--\ref{lem:tail}). At $\eps=0$ the weights are $\binom dj$, with sum $2^d$; they concentrate
at $j=d/2$, where $F=\ln2$, so $\thit^+/\thit\to e^{c\ln2}=2^c$.

\begin{theorem}\label{thm:main}
For the simple random walk on $Q_d$ and $\eps=c/d$ with $c>0$ fixed,
\[
\lim_{d\to\infty}\frac{\thit^+(c/d)}{\thit}=2^{c}.
\]
Consequently $d\,\rhit(C)\to\log_2 C$ for every $C>1$.
\end{theorem}

\begin{theorem}\label{thm:lazy}
For the lazy walk on $Q_d$ (holding probability $1/2$, each neighbour $1/(2d)$; the adversary also perturbs the holding
probability by a factor in $[1-\eps,1+\eps]$, and $\thit$, $\thit^+$ refer to the lazy walk) and $\eps=c/d$ with $c>0$
fixed, the ratio $\thit^+(c/d)/\thit$ tends to $e^{c}$, and $d\,\rhit(C)\to\ln C$ for every $C>1$.
\end{theorem}

In both cases $\rhit$ is of order $1/\Delta$, so the hypercube satisfies the conjecture $\rhit=\Omega(1/\Delta)$
of~\cite{Zuc26}, with explicit constants. The extremal adversary is stronger than the tilted adversary: it multiplies $\thit$ by $e^{\eps d\ln 2}$ against about $e^{\eps d/2}$, so the limit
$\log_2C$ of $d\,\rhit(C)$ lies below the bound $2\ln C$ that the tilted adversary gives.

Theorem~\ref{thm:main} can also be stated through stationary probabilities. The walk on $Q_d$ gives the target $0$
stationary probability $\pi(0)=2^{-d}$. Under the extremal adversary the walk is again a Markov
chain, and its stationary probability of $0$ is $\pi^+(0)=1/W$. By $T_1=W-1$ and Lemma~\ref{lem:tail},
$\thit\sim2^d=1/\pi(0)$ and $\thit^+\sim W=1/\pi^+(0)$, so Theorem~\ref{thm:main} says
$\pi^+(0)=\pi(0)^{1+\eps+o(\eps)}$ at $\eps=c/d$. Azar, Broder,
Karlin, Linial and Phillips~\cite{ABKLP96}, in a related model where a controller chooses the step with probability
$\eps$, showed that on bounded-degree graphs the controller can raise the stationary probability of a vertex $v$ to
$\pi(v)^{1-\Omega(\eps)}$, and Haslegrave, Sauerwald and Sylvester~\cite{HSS22} extended this to further
settings. The hypercube has unbounded degree, the adversary here lowers the probability of the target,
and the exponent is exact.

\section{Reduction to a birth--death chain}\label{sec:reduction}

Fix the target $0\in Q_d$ and write $|v|$ for the Hamming weight. For $0\le k\le d$ put
\[
q_k=\max\Bigl(\frac{k(1-\eps)}{d},\ 1-\frac{(d-k)(1+\eps)}{d}\Bigr),\qquad p_k=1-q_k .
\]
The first term is the larger exactly when $k\le d/2$. Let $V_k$ be the expected time for the birth--death chain on
$\{0,\dots,d\}$ that steps from $k$ to $k-1$ with probability $q_k$ and to $k+1$ with probability $p_k$ to reach $0$
from $k$, and $T_k=V_k-V_{k-1}$ the expected time to go from $k$ to $k-1$. Then $T_d=1/q_d=1$,
$T_k=(1+p_kT_{k+1})/q_k$, and in particular every $T_k$ is positive, so $V$ is increasing.

\begin{lemma}\label{lem:reduction}
For every standard $\eps$-adversary and every $v$, $\E_vH_0\le V_{|v|}$, with equality for the extremal adversary,
which at a vertex of weight $k$ gives each of the $k$ neighbours of weight $k-1$ probability $q_k/k$ and each of the $d-k$
neighbours of weight $k+1$ probability $p_k/(d-k)$. Hence $\thit^+(\eps)=V_d$.
\end{lemma}

\begin{proof}
The extremal adversary is admissible. For $k\le d/2$, $q_k/k=(1-\eps)/d$ and $p_k/(d-k)=(1+\eps k/(d-k))/d$, which lies
in $[1/d,(1+\eps)/d]$ because $k\le d-k$. For $k>d/2$, $p_k/(d-k)=(1+\eps)/d$ and $q_k/k=(1-\eps(d-k)/k)/d$, which
lies in $[(1-\eps)/d,1/d]$. Under it $|X_s|$ is the birth--death chain above, so $\E_vH_0=V_{|v|}$.

Now let the adversary be arbitrary, let $X_s$ be the walk and $\mathcal F_s$ its history up to time $s$, and put
$U(v)=V_{|v|}$. At a vertex $v\ne0$ of weight $k$, any admissible step gives total probability $m\ge q_k$ to the neighbours of weight $k-1$, since each of them receives at least
$(1-\eps)/d$ and the others together at most $(d-k)(1+\eps)/d$. As $V_{k-1}<V_{k+1}$,
\[
\E\bigl[U(X_{s+1})\,\big|\,\mathcal F_s\bigr]=mV_{k-1}+(1-m)V_{k+1}\le q_kV_{k-1}+p_kV_{k+1}=U(X_s)-1
\]
on $\{X_s=v\}$, where the last equality is the first-step equation of the birth--death chain (for $k=d$ both sides
equal $V_{d-1}$). So $U(X_{s\wedge H_0})+s\wedge H_0$ is a supermartingale, and $\E_v[U(X_{s\wedge H_0})]+\E_v[s\wedge
H_0]\le U(v)$ for every $s$. Since $U\ge0$, letting $s\to\infty$ gives $\E_vH_0\le U(v)$ by monotone convergence.
Finally $\thit^+=\max_{u,v}h^+(u,v)$, the hypercube is vertex-transitive, and $V$ is increasing, so
$\thit^+=V_d$.
\end{proof}

On every graph, history does not help the adversary~\cite[Section~1.4]{Zuc26}; on $Q_d$ this also follows from the
proof above.

For the lazy walk of Theorem~\ref{thm:lazy} put $\tilde q_k=k(1-\eps)/(2d)$ and $\tilde p_k=(d-k)(1+\eps)/(2d)$, and let
$\tilde V_k$ be the expected time for the birth--death chain on $\{0,\dots,d\}$ that steps from $k$ to $k-1$ with
probability $\tilde q_k$, to $k+1$ with probability $\tilde p_k$ and otherwise stays at $k$, to reach $0$ from $k$. Put
$\tilde T_k=\tilde V_k-\tilde V_{k-1}$. The holding term cancels in the first-step equation, so
$\tilde T_d=1/\tilde q_d$ and $\tilde T_k=(1+\tilde p_k\tilde T_{k+1})/\tilde q_k$, and $\tilde V$ is increasing.

\begin{lemma}\label{lem:lazy}
For every adversary of Theorem~\ref{thm:lazy} and every $v$, $\E_vH_0\le\tilde V_{|v|}$, with equality for the
adversary that at a vertex of weight $k$ gives each neighbour of weight $k-1$ probability $(1-\eps)/(2d)$ and each
neighbour of weight $k+1$ probability $(1+\eps)/(2d)$. Hence $\thit^+(\eps)=\tilde V_d$ for the lazy walk.
\end{lemma}

\begin{proof}
The holding probability of that adversary is $1-\tilde q_k-\tilde p_k=\tfrac12+\eps(2k-d)/(2d)$, which lies in
$[(1-\eps)/2,(1+\eps)/2]$, so it is admissible, and under it $|X_s|$ is the chain above. For an arbitrary adversary
let $m_-$, $m_0$, $m_+$ be the masses a step from a vertex of weight $k$ puts on the neighbours of weight $k-1$, on
holding and on the neighbours of weight $k+1$. Then $m_-\ge\tilde q_k$ and $m_+\le\tilde p_k$, and with
$U(v)=\tilde V_{|v|}$, since $\tilde V_{k-1}<\tilde V_k<\tilde V_{k+1}$,
\[
\E\bigl[U(X_{s+1})\,\big|\,\mathcal F_s\bigr]=m_-\tilde V_{k-1}+m_0\tilde V_k+m_+\tilde V_{k+1}
\le\tilde q_k\tilde V_{k-1}+(1-\tilde q_k-\tilde p_k)\tilde V_k+\tilde p_k\tilde V_{k+1}=U(X_s)-1 .
\]
The rest of the proof of Lemma~\ref{lem:reduction} applies unchanged.
\end{proof}

\section{Proof of the theorems}\label{sec:proof}

Put $w_0=1$, $w_j=\prod_{i<j}p_i/q_{i+1}$ and $W=\sum_{j=0}^dw_j$. Since $q_{k+1}w_{k+1}=p_kw_k$, the recursion for
$T_k$ is solved by the classical formula for birth--death chains~\cite[Section~2.5]{LPW17}
\begin{equation}\label{eq:Tk}
T_k=\frac{1}{q_kw_k}\sum_{j\ge k}w_j\qquad(1\le k\le d).
\end{equation}
For the simple walk $q_0=0$, so $p_0=1=q_1w_1$ and $T_1=W-1$. For $\eps=0$, $w_j=\binom dj$ and $W=2^d$.

\begin{lemma}\label{lem:weights}
Let $\eps=c/d$ with $c\ge0$ fixed. Uniformly in $0\le j\le d$, with implied constants depending only on $c$,
\[
\ln\Bigl(w_j\Big/\binom dj\Bigr)=c\,F(j/d)+O(1/d),\qquad
F(x)=\begin{cases}-\ln(1-x),&0\le x\le\tfrac12,\\ \ln(4x),&\tfrac12\le x\le1.\end{cases}
\]
\end{lemma}

\begin{proof}
Compare the factor $p_i/q_{i+1}$ with its value $(d-i)/(i+1)$ at $\eps=0$. If $i+1\le d/2$ the ratio is
$(1+\eps i/(d-i))/(1-\eps)$, whose logarithm is $\eps d/(d-i)+O(\eps^2)$. If $i\ge d/2$ it is
$(1+\eps)(i+1)/(i+1-\eps(d-i-1))$, whose logarithm is $\eps d/(i+1)+O(\eps^2)$. At most one index lies between the
two cases, contributing $O(\eps)$. With $\eps=c/d$ the logarithm of the factor is $(c/d)g(i/d)+O(1/d^2)$, where
$g(y)=1/(1-y)$ for $y<1/2$ and $g(y)=1/y$ for $y\ge1/2$. Summing over $i<j$ and comparing the Riemann sum of the
bounded monotone pieces of $g$ with $F(x)=\int_0^xg$ gives the claim.
\end{proof}

In particular $K^{-1}\binom dj\le w_j\le K\binom dj$ for all $j$, with $K=e^{c\ln4+O(1/d)}$ bounded.

\begin{lemma}\label{lem:tail}
For $\eps=c/d$ with $c\ge0$ fixed, $\sum_{k\ge2}T_k=O(W/d+d\log d)$, the implied constant depending only on $c$.
\end{lemma}

\begin{proof}
For $2\le k\le d/2$ every factor in $w_k$ is at least its value at $\eps=0$, so $w_k\ge\binom dk$, and
$q_k\ge k(1-\eps)/d$. By~\eqref{eq:Tk}, $T_k\le Wd/(k(1-\eps)\binom dk)$, and
$\sum_{k=2}^{d/2}d/(k\binom dk)=O(1/d)$. For $k>d/2$, $q_k\ge(1-\eps)/2$ and, by the bounds after
Lemma~\ref{lem:weights},
\[
T_k\le\frac{2K^2}{1-\eps}\sum_{j\ge k}\frac{\binom dj}{\binom dk}\le\frac{2K^2}{1-\eps}\cdot\frac{k+1}{2k+1-d},
\]
since consecutive ratios $\binom d{j+1}/\binom dj=(d-j)/(j+1)$ are at most $(d-k)/(k+1)<1$ for $j\ge k$. Summing over
$d/2<k\le d$ gives $O(d\log d)$.
\end{proof}

\begin{proof}[Proof of Theorem~\ref{thm:main}]
By Lemma~\ref{lem:reduction}, $\thit^+(\eps)=V_d=\sum_kT_k=W-1+O(W/d+d\log d)$ by Lemma~\ref{lem:tail}, and the same
holds at $\eps=0$ with $W=2^d$. By Lemma~\ref{lem:weights},
\[
\frac{W}{2^d}=\E\bigl[e^{cF(B/d)+O(1/d)}\bigr],\qquad B\sim\mathrm{Bin}(d,\tfrac12).
\]
$F$ is bounded and continuous and $B/d\to\frac12$ in probability, so the right side tends to $e^{cF(1/2)}=e^{c\ln2}=
2^c$, and $\thit^+(c/d)/\thit\to2^c$.

For the second statement, $\thit^+(\eps)$ is nondecreasing in $\eps$. If $c<\log_2C$ then $\thit^+(c/d)\le C\thit$ for
all large $d$, so $\liminf_d d\,\rhit(C)\ge c$; if $c>\log_2C$ then $\thit^+(c/d)>C\thit$ for all large $d$, so
$\limsup_d d\,\rhit(C)\le c$.
\end{proof}

\begin{proof}[Proof of Theorem~\ref{thm:lazy}]
By Lemma~\ref{lem:lazy}, $\thit^+(\eps)=\tilde V_d=\sum_k\tilde T_k$. Here $\tilde p_i/\tilde q_{i+1}=r\,(d-i)/(i+1)$
with $r=(1+\eps)/(1-\eps)$, so the weights $\tilde w_j=\prod_{i<j}\tilde p_i/\tilde q_{i+1}$ are $\binom djr^j$ and
$\tilde W=\sum_j\tilde w_j=(1+r)^d=(2/(1-\eps))^d$ exactly. The recursion for $\tilde T_k$ is the one for $T_k$, so
\eqref{eq:Tk} holds with tildes; $\tilde q_1\tilde w_1=\tilde p_0=(1+\eps)/2$ gives $\tilde T_1=2(\tilde W-1)/(1+\eps)$.
The proof of Lemma~\ref{lem:tail} applies with $\tilde K=r^d=e^{2c+O(1/d)}$ in place of $K$; there
$\tilde q_k=k(1-\eps)/(2d)$ is half the value used in that proof, so $\tilde q_k\ge(1-\eps)/4$ for $k>d/2$ and the
constants double, and $\sum_{k\ge2}\tilde T_k=O(\tilde W/d+d\log d)$. At $\eps=0$, $\tilde T_1=2(2^d-1)$. Therefore
\[
\frac{\thit^+(c/d)}{\thit}=\frac{(1-c/d)^{-d}}{1+c/d}\,(1+o(1))\to e^{c},
\]
and $d\,\rhit(C)\to\ln C$ as before.
\end{proof}

\begin{remark}
The limits are the same for the interval $[e^{-\eps}P,e^{\eps}P]$, which~\cite[Section~1.1]{Zuc26} also uses: with
$\eps=c/d$ the two intervals differ by $O(\eps^2)$ in relative terms, which changes each $\ln w_j$ by $O(d\eps^2)=O(1/d)$; the limits in both theorems are unchanged.
\end{remark}

\section{Numerical values}\label{sec:checks}

The table shows $d\,\rhit(2)$ within $0.002$ of $1$ from $d=32$ on, while $d\,\rhit(4)$ stays below $2$ at every $d$
listed and is still about $0.005$ short at $d=512$. Lemma~\ref{lem:reduction} reduces $\thit^+$ to a linear recursion, which the program \texttt{rho\_hit.py} in the
repository~\cite{Code} solves numerically, finding $\rhit(C)$ by bisection. For the simple walk:
\begin{center}
\begin{tabular}{rcccc}
$d$ & $\thit/2^d$ & $\rhit(2)$ & $d\,\rhit(2)$ & $d\,\rhit(4)$\\\hline
8 & 1.2190 & 0.13142 & 1.0514 & 1.9334\\
16 & 1.0798 & 0.06286 & 1.0057 & 1.9249\\
32 & 1.0347 & 0.03120 & 0.9985 & 1.9493\\
64 & 1.0164 & 0.01560 & 0.9981 & 1.9702\\
128 & 1.0080 & 0.00780 & 0.9987 & 1.9835\\
256 & 1.0040 & 0.00390 & 0.9992 & 1.9912\\
512 & 1.0020 & 0.00195 & 0.9996 & 1.9954
\end{tabular}
\end{center}
For the lazy walk (\texttt{lazy\_check.py}) $d\,\rhit(2)=0.6940$ and $d\,\rhit(4)=1.3871$ at $d=512$, against $\ln2=
0.6931$ and $\ln4=1.3863$.

\section*{Acknowledgements and funding}
This work was done by AI agents, Claude (Anthropic), connected to Math Brain and working inside the Brain framework, a
shared, persistent memory in which the agents record claims, proofs, failed attempts and methods, retrieve one
another's results, coordinate who works on what, and audit each other's work. For more on the Brain and Math Brain,
see the Math Brain paper~\cite{MathBrainPaper}, the Math Brain repository~\cite{MathBrainRepo} and the Brain framework
repository~\cite{BrainFramework}.

The research presented in this paper was supported by the European Research Council (ERC) under the European Union's
Horizon Europe research and innovation programme (grant agreement No.~101041711), by the Simons Foundation as part of
the Collaboration on the Mathematical and Scientific Foundations of Deep Learning, by Heights Labs, by Convex Nexus
Capital, by the Israel Science Foundation (grant number 2258/19), and by the Israel Science Foundation (ISF Grant
4101/25). Computational resources were provided in part by Google Cloud through its education credits programme
(credit 530550103).

\section*{Declaration of competing interest}
The author holds equity in Heights Labs and in Convex Nexus Capital and has a management role at Convex Nexus Capital.
These entities are included in the funding declaration above.

\section*{Data availability}
The programs described in Section~\ref{sec:checks} are in the repository~\cite{Code}.

\section*{Declaration of generative AI and AI-assisted technologies}
This paper is part of a broader research project on AI proving problems. Claude agents (Anthropic) proved
Theorems~\ref{thm:main} and~\ref{thm:lazy}, wrote the programs of Section~\ref{sec:checks} and prepared the manuscript,
and a separate Claude model reviewed the proofs. The model, the definitions and the question are due to
Zuckerman~\cite{Zuc26}, and the upper bound $\rhit=O(1/d)$ on $Q_d$ follows from the adversary constructed in the proof of his Theorem~4.5. The author directed
the work and is responsible for the paper.

\begin{thebibliography}{9}
\bibitem{ABKLP96} Y.~Azar, A.~Z.~Broder, A.~R.~Karlin, N.~Linial and S.~J.~Phillips, \emph{Biased random walks},
Combinatorica \textbf{16} (1996), 1--18, \doi{10.1007/BF01300124}.
\bibitem{HSS22} J.~Haslegrave, T.~Sauerwald and J.~Sylvester, \emph{Time dependent biased random walks}, ACM Trans.\ Algorithms \textbf{18} (2022).
\bibitem{LPW17} D.~A.~Levin and Y.~Peres, \emph{Markov chains and mixing times}, 2nd ed., with contributions by
E.~L.~Wilmer, American Mathematical Society, Providence, RI, 2017, \doi{10.1090/mbk/107}.
\bibitem{BrainFramework} Y.~Shmalo, \emph{Brain Framework}, public software repository, 2026,
\url{https://github.com/yspennstate/brain-framework-public/tree/5629e644cdd510173ae97e9aa92642dabd22dc57}.
\bibitem{Code} Y.~Shmalo, \emph{Hitting-time robustness of the hypercube: source code and paper}, 2026,
\url{https://github.com/yspennstate/hypercube-hitting-robustness}.
\bibitem{MathBrainRepo} Y.~Shmalo, \emph{Math Brain: public snapshot, code and tests}, 2026,
\url{https://github.com/yspennstate/math-brain-public/tree/8c9a6bd478c4ec4082347151b7ae16db50c917dd}.
\bibitem{MathBrainPaper} Y.~Shmalo, \emph{Math Brain: shared, auditable memory for multi-agent mathematical research},
2026, \doi{10.5281/zenodo.24150837}.
\bibitem{Zuc26} D.~Zuckerman, \emph{Markov chain robustness}, in: 17th Innovations in Theoretical Computer Science
Conference (ITCS 2026), LIPIcs~362, Schloss Dagstuhl--Leibniz-Zentrum f\"ur Informatik, 2026, Paper~118,
\doi{10.4230/LIPIcs.ITCS.2026.118}.
\end{thebibliography}

\end{document}
